\documentclass{ximera}

\input{../../preamble.tex}

\title{Chaining Functions}

\begin{document}

\begin{abstract}
composition
\end{abstract}
\maketitle





Using two given two functions, we can create a third function, called the \textbf{composition}.

If $f$ and $g$ are the names of the two component funcitons, then $f \circ g$ is the symbol for the compositoin of $f$ and $g$.





\begin{idea} \textbf{\textcolor{green!50!black}{Composition}}


Let $F$ and $G$ be two functions.

The composition of $F$ and $G$ is another function called the composition of $F$ and $G$.  A circle between the function names denotes a composition.

\[
F \circ G
\]


\[
(F \circ G)(a) = F(G(a))
\]


The composition applies $G$ to an element of the domain of $G$ and then uses the resulting function value from $G$ as a domain number of $F$ obtaining a function value of $F$.

This ultimate value of $F$ from this process is the value of $F \circ G$.

\end{idea}



Composition is our way of representing an ordered application of processes.














\begin{onlineOnly}
\begin{center}
\textbf{\textcolor{green!50!black}{ooooo-=-=-=-ooOoo-=-=-=-ooooo}} \\

more examples can be found by following this link\\ \link[More Examples of Formulas]{https://ximera.osu.edu/csccmathematics/precalculus1/precalculus1/formulas/examples/exampleList}

\end{center}
\end{onlineOnly}





\end{document}
